\subsection*{小结}

\begin{table}[htbp]\centering\small
\begin{tabularx}{\linewidth}{@{}lXX@{}}
\toprule
\textbf{工具} & \textbf{作用阶段} & \textbf{开销或行为} \\
\midrule
\co@print!@ / \co@println!@ & 编译时 & 不增加运行时开销 \\
\co@assert(expr, "msg")@ & 运行时 & 开销较低，条件不满足时终止执行 \\
\co@print@ / \co@println@ & 运行时 & 打印可能串行执行，应限制输出范围 \\
\co@-rtc=high@ & 运行时 & 增加边界等检查 \\
\co@-v@ / \co@--verbose@ & 编译过程 & 显示外部命令，不增加程序运行时开销 \\
\co@-g@、调试器与 RTTI & 调试过程 & 提供符号和类型信息，性能取决于构建配置 \\
\bottomrule
\end{tabularx}
\end{table}

这些工具观察的是同一程序的不同阶段。编译器输出说明 Croqtile 推导出了哪些形状和类型，运行时检查则暴露某次调用中的实际取值。构建失败时，详细模式能显示失败的外部命令，便于单独复现。构建成功与结果正确，需要分别确认。

小例子验证通过后，再把同样的方法用于需要修改的核函数。保留参考结果，每次调整程序的一处内容，先比较输出，再测量性能。这样出现回退时，才能把问题对应到具体修改。
